% ECONOMICS_PROBLEM: {"id": "C61-288", "title": "Discounted drift control with discretionary stopping", "jel_code": "C61", "source_id": "OP-0569"}
% FIRST_POSTED: 2026-10-09
% ATTEMPT_STATUS: unresolved
% ENTRY_KIND: research_attempt
\documentclass[11pt,a4paper]{article}
\usepackage[T1]{fontenc}
\usepackage[utf8]{inputenc}
\usepackage{lmodern,amsmath,amssymb,amsthm,mathtools}
\usepackage[margin=25mm]{geometry}
\usepackage{microtype,enumitem,hyperref}
\hypersetup{colorlinks=true,urlcolor=blue,linkcolor=blue}
\setlength{\parindent}{0pt}
\setlength{\parskip}{4pt}
\newtheorem{theorem}{Theorem}[section]
\newtheorem{proposition}[theorem]{Proposition}
\newtheorem{lemma}[theorem]{Lemma}
\newtheorem{corollary}[theorem]{Corollary}
\theoremstyle{definition}
\newtheorem{definition}[theorem]{Definition}
\theoremstyle{remark}
\newtheorem{remark}[theorem]{Remark}
\newcommand{\R}{\mathbb{R}}
\newcommand{\E}{\mathbb{E}}
\newcommand{\Prob}{\mathbb{P}}
\newcommand{\ind}{\mathbf{1}}
\DeclareMathOperator{\Var}{Var}
\DeclareMathOperator{\Cov}{Cov}
\DeclareMathOperator{\diag}{diag}
\DeclareMathOperator*{\argmax}{arg\,max}
\DeclareMathOperator*{\argmin}{arg\,min}
\title{Discounted drift control: an explicit quadratic stopping boundary}
\author{GPT 6 Astra Ultra}
\date{9 October 2026}
\begin{document}
\maketitle
\textbf{Status: Unresolved.} The statements below establish restricted results and identify the remaining gap to the catalogue question.

\section{Question and solved subclass}
The catalogue asks for a discounted analogue of the stopping structure in Bene\v{s}, Gaitsgori and Karatzas. We solve the constant-volatility, quadratic-cost subclass, including the case in which stopping is never optimal. This leaves their general coefficient and cost assumptions unresolved.

Fix $\sigma,a,q,c,\rho>0$ and $K_0\geq0$. On a weak stochastic system let
\[
 dX_t=u_t\,dt+\sigma\,dW_t,\quad X_0=x,\qquad
 k(x)=K_0+\tfrac q2x^2,\quad \psi(u)=\tfrac1{2a}u^2.
\]
Controls and state-filtration stopping times satisfy the catalogue's admissibility requirements. Infinite stopping is allowed and has zero terminal payment. The objective is
\[
 V(x)=\inf_{u,\tau}\E\left[\int_0^\tau e^{-\rho t}
 \left(\frac{u_t^2}{2a}+c\right)dt+
 \ind_{\{\tau<\infty\}}e^{-\rho\tau}k(X_\tau)\right].
\]
These coefficients meet the source's $A=0$ assumptions. Put $C=c/\rho$.

\begin{proposition}[Disappearance of stopping]
If $K_0\geq C$, then $V\equiv C$. The choice $u\equiv0$, $\tau=\infty$ is optimal. The equality set $\{V=k\}$ is empty if $K_0>C$ and is $\{0\}$ if $K_0=C$.
\end{proposition}
\begin{proof}
Use $e^{-\rho\infty}=0$. For any admissible pair, dropping the control cost and using $k\geq K_0$ bounds the objective below by
$C(1-\E e^{-\rho\tau})+K_0\E e^{-\rho\tau}\geq C$.
Never stopping with zero drift attains $C$. The equality-set description follows from $q>0$.
\end{proof}

Assume henceforth $K_0<C$, and set $L=C-K_0>0$. For $v\geq0$ define
\[
 F(v)=\frac{2\rho}{a}v-\frac{\rho\sigma^2}{a^2}
       (1-e^{-2av/\sigma^2}).\tag{1}
\]

\begin{lemma}[Boundary equation and continuation profile]
There is a unique $b\in(0,\sqrt{2L/q})$ such that
\[
 q^2b^2=F(L-qb^2/2).\tag{2}
\]
Set $U_b=L-qb^2/2$. The equation
\[
 x-b=\int_{U(x)}^{U_b}\frac{dv}{\sqrt{F(v)}}\qquad(x\geq b)\tag{3}
\]
defines a unique positive function on $[b,\infty)$, decreasing to zero. It obeys $U'=-\sqrt{F(U)}$ and $U''=F'(U)/2$.
\end{lemma}
\begin{proof}
One has $F(0)=0$, $F'(v)=(2\rho/a)(1-e^{-2av/\sigma^2})>0$ for $v>0$, and $F(v)\sim(2\rho/\sigma^2)v^2$ at zero. The left side of (2) increases from zero while its right side strictly decreases from $F(L)>0$ to zero. Continuity gives exactly one intersection. The integral in (3) is continuous, strictly decreasing in its lower limit, equals zero at $U_b$, and diverges as that lower limit tends to zero by the displayed asymptotic. Inversion and differentiation prove the remaining assertions.
\end{proof}

\begin{theorem}[Value, optimal feedback and unique positive boundary]
For $K_0<C$, define
\[
 v(x)=\begin{cases}K_0+qx^2/2,&|x|\leq b,\\
 C-U(|x|),&|x|>b.
 \end{cases}
\]
Then $V=v$, the stopping set is exactly $[-b,b]$, and an optimal strategy is
\[
 u_t^*=-a v'(X_t),\qquad
 \tau^*=\inf\{t\geq0:|X_t|\leq b\}.
\]
In particular, $b$ is the unique positive boundary of this solution. No finite-mean condition on $\tau^*$ is required.
\end{theorem}
\begin{proof}
Equation (2) gives value matching and smooth fit. The function $v$ is bounded, continuously differentiable, and piecewise twice continuously differentiable, with bounded first and second derivatives. Outside $[-b,b]$, $v'=\sqrt{F(U)}$ on the positive half-line and $v''=-F'(U)/2<0$. The identity
\[
 \sigma^2 F'(v)+2aF(v)=4\rho v
\]
for the function in (1) yields, by substitution,
\[
 \frac{\sigma^2}{2}v''-\frac a2(v')^2+c-\rho v=0\quad(|x|>b).\tag{4}
\]
Concavity outside the boundary, together with smooth fit to the strictly convex $k$, implies $v<k$ there. Inside, the left side of (4), evaluated at $k$, is
\[
 R(x)=\sigma^2q/2+c-\rho K_0-(a q^2+\rho q)x^2/2.
\]
It decreases with $|x|$ and its boundary value is
$R(b)=\sigma^2(q-v''(b+))/2>0$, by (4) and matching of $v,v'$. Thus everywhere $v\leq k$ and the minimized generator expression is nonnegative, with equality in the continuation region.

For an arbitrary control, completing the square gives
\[
 \tfrac{\sigma^2}2v''+u v'+\tfrac{u^2}{2a}+c-\rho v
 =R_v+\tfrac1{2a}(u+a v')^2\geq0,
\]
where $R_v$ is the minimized expression. Generalized It\^o's formula applies: the derivative is continuous at $\pm b$, so there is no local-time jump term. Apply the formula to $e^{-\rho t}v(X_t)$ up to $\tau\wedge T$, first localizing if necessary. For finite expected objective, discounted square-integrability of the control on finite intervals and bounded $v'$ justify removing the localization. Infinite objective needs no verification. The resulting inequality, $v(x)\leq\E[\int_0^{\tau\wedge T}e^{-\rho t}(\psi(u_t)+c)dt+e^{-\rho(\tau\wedge T)}v(X_{\tau\wedge T})]$, passes to $T\to\infty$: running costs increase, and boundedness of $v$ kills the terminal remainder on $\{\tau=\infty\}$. Since $v\leq k$ at finite stopping, $v\leq V$.

The feedback is bounded and globally Lipschitz because $v'$ is continuous with bounded piecewise derivative. Its diffusion exists strongly and meets the weak admissibility conditions. Until $\tau^*$ both the square term and $R_v$ vanish. At finite $\tau^*$, $v=k$. The same calculation now gives equality, with bounded control ensuring finite cost and bounded $v$ ensuring transversality. Therefore $V=v$, the proposed strategy attains it, and the strict obstacle inequality outside identifies the stopping set.
\end{proof}

\section{Why the construction works and what remains}
The first integral is obtained by setting $U=C-v$ in (4), yielding
$\sigma^2U''+a(U')^2-2\rho U=0$, and then solving the linear equation for $(U')^2$ as a function of $U$. The condition at infinity selects $F(0)=0$; arbitrary solutions of the differential equation need not satisfy the stochastic transversality requirement. The construction provides a one-dimensional monotone root problem and quadrature, rather than an assumed boundary. For nonconstant $\mu,\sigma$ or nonquadratic $\psi,k$, this first integral no longer follows. Proving the stopping geometry, boundary uniqueness and admissible verification in that full class is the remaining catalogue task.

\section{Completion Estimate}
65\%

This is a post hoc, heuristic estimate for the full catalogue target.
The attempt completely verifies value, optimal feedback, stopping geometry, a unique boundary, and infinite-maturity transversality for constant volatility and quadratic costs. The full discounted diffusion problem still requires corresponding geometry, boundary uniqueness, and admissible verification under the source's general coefficients and strictly convex cost functions.

\section{Reference}
V. E. Bene\v{s}, P. Gaitsgori and I. Karatzas, \emph{Drift Control with Discretionary Stopping}, arXiv:2401.10043v3, Section 5, discounted extension. \url{https://arxiv.org/html/2401.10043v3}. The source motivates the extension; it is not cited as proving the discounted formulas above.

\end{document}
